\documentclass[11pt]{article}
\usepackage[T1]{fontenc}
\usepackage{lmodern}
\usepackage[margin=1in]{geometry}
\usepackage{amsmath,amssymb,amsthm,mathtools}
\usepackage{microtype,booktabs,longtable,array,enumitem,xurl}
\usepackage[dvipsnames]{xcolor}
\usepackage[colorlinks=true,linkcolor=MidnightBlue,citecolor=MidnightBlue,urlcolor=MidnightBlue]{hyperref}
\usepackage{bookmark}
\usepackage{fancyhdr}
\pagestyle{fancy}
\fancyhf{}
\fancyhead[L]{\small Rational Herglotz companion values}
\fancyhead[R]{\small ProveIt research continuation}
\fancyfoot[C]{\thepage}
\setlength{\headheight}{14pt}
\setlength{\emergencystretch}{2em}
\setcounter{tocdepth}{2}
\numberwithin{equation}{section}
\newtheorem{theorem}{Theorem}[section]
\newtheorem{proposition}[theorem]{Proposition}
\newtheorem{lemma}[theorem]{Lemma}
\newtheorem{corollary}[theorem]{Corollary}
\theoremstyle{definition}
\newtheorem{definition}[theorem]{Definition}
\newtheorem{example}[theorem]{Example}
\theoremstyle{remark}
\newtheorem{remark}[theorem]{Remark}
\newtheorem{question}[theorem]{Research question}
\newcommand{\Q}{\mathbb Q}
\newcommand{\R}{\mathbb R}
\newcommand{\C}{\mathbb C}
\newcommand{\Z}{\mathbb Z}
\newcommand{\N}{\mathbb N}
\newcommand{\A}{\mathcal A}
\newcommand{\B}{\mathcal B}
\newcommand{\Pcal}{\mathcal P}
\newcommand{\Rog}{\mathcal L}
\newcommand{\cl}[1]{\langle #1\rangle}
\newcommand{\ee}{\mathbf e}
\newcommand{\dd}{\,\mathrm d}
\newcommand{\inv}{^{-1}}
\DeclareMathOperator{\Li}{Li}
\DeclareMathOperator{\Gal}{Gal}
\DeclareMathOperator{\Span}{Span}
\DeclareMathOperator{\Norm}{Norm}
\newcommand{\snapshot}{\texttt{cc34f73596336f2466d9754cb0f3635bd2bedade}}
\hypersetup{pdftitle={Rational Herglotz Companion Values: Dyadic Descent, Complete Formal Classification, and Explicit Identities},pdfauthor={OpenAI-assisted research continuation for the ProveIt project},pdfsubject={Dilogarithms, cyclotomic exterior symbols, Herglotz function, exact certificates}}
\begin{document}
\begin{titlepage}
\thispagestyle{empty}
\vspace*{0.45in}
{\sffamily\small PROVEIT \quad / \quad ANALYSIS AND POLYLOGARITHMS\par}
\vspace{0.55in}
{\LARGE\bfseries Rational Herglotz\par Companion Values\par}
\vspace{0.23in}
{\Large Dyadic descent, complete formal classification,\par and explicit dilogarithm identities\par}
\vspace{0.45in}
{\large Research continuation prepared for the ProveIt project\par}
\vspace{0.10in}
{\normalsize OpenAI-assisted mathematical research note\par October 2026\par}
\vspace{0.5in}
\begin{abstract}
\noindent
For the Herglotz companion
$J(x)=\int_0^1\log(1+t^x)/(1+t)\dd t$, we resolve the manuscript's
all-rational \emph{formal} reduction problem. At a reduced mixed-parity
ratio, with even entry $e$ and odd entry $o$, reduction to logarithms is
equivalent to $e^2\equiv\pm1\pmod o$ and $o^2\equiv1\pmod{2e}$.
Allowing rational dilogarithms adds precisely six rational arguments:
$1/3,3,1/5,5,3/5,5/3$. The proof combines a conductor-doubling identity
with rigidity of the new dyadic character layer. It uses the existing
all-conductor symbol theorem, whose proof is recalled.

We give a complete even-coefficient quadratic-recurrence parametrization
of the logarithmic cases and an all-finite-order height-counting expansion.
Analytically, a Hecke relation and the Rogers five-term identity yield an
explicit log-sine formula for every $J(n/(n+1))$, and a compact evaluation
of $J(3/5)$. A prime-seven exterior-logarithm witness disproves a
finite-field three-term assertion in the reviewed Radchenko--Zagier
preprint, without affecting the direct boundary formula used here.
Exact finite tests and independent high-precision integral comparisons
accompany the proofs. Formal non-reducibility is not numerical
transcendence; no period-injectivity or proof-assistant claim is made.
\end{abstract}
\vfill
\noindent\small
\textbf{Source snapshot.} \snapshot\par
\noindent\small
\textbf{Target.} \nolinkurl{Analysis/Polylogarithms/docs/manuscript}\par
\noindent\small
\textbf{Package.} LaTeX source, compiled PDF, exact arithmetic checks,
independent numerical diagnostics, theorem ledger and integration notes.
\end{titlepage}
\tableofcontents
\clearpage
\section{The question being continued}
\label{sec:context}

The current ProveIt manuscript treats special polylogarithmic values,
cyclotomic relations, exact certificates and Herglotz arithmetic in one
framework. Its Herglotz chapter proves a complete formal classification
of $F(p/q)-F(1)$ and of $J(2/q)$. Its discovery chapter explicitly leaves
the general $J(p/q)$ symbol problem as a further direction
\cite{Repo}. That is the problem addressed here. The Gaussian reductions,
the $S_4$ certificate and the other completed continuations are not
reclaimed as new work.

For $x>0$, use the standard normalization
\begin{equation}\label{eq:def-FJ}
 F(x)=\sum_{n\ge1}\frac{\psi(nx)-\log(nx)}{n},
 \qquad J(x)=\int_0^1\frac{\log(1+t^x)}{1+t}\dd t.
\end{equation}
Radchenko--Zagier establish the finite cyclotomic dilogarithm formula for
rational $F$ and the identity
\begin{equation}\label{eq:J-F}
 J(x)=F(2x)-2F(x)+F(x/2)+\frac{\pi^2}{12x}.
\end{equation}
Their modified Herglotz function has a Hecke relation that will supply
explicit evaluations in Section~\ref{sec:identities} \cite{RZ}.
The broader two-parameter Herglotz--Zagier--Novikov function and its
specializations are studied by Choie--Kumar \cite{CK}.

There are two different questions. An analytic evaluation is an equality
of numbers with specified branches. A formal reduction asks whether the
finite dilogarithm expression admits reduction by the rationalized
five-term calculus. Section~\ref{sec:symbols} defines the latter exactly.
The classification below is complete for that calculus. It does not
classify every conceivable numerical equality among these periods.

\subsection{Results and their dependencies}

The new algebraic step is not another numerical rank computation. At an
odd conductor, doubling the root set gives an exact three-term
\emph{conductor} identity for exterior symbols. At an even conductor,
a symbol in the new dyadic layer cannot be absorbed by symbols from
the preceding field unless it is zero. These facts reduce the boundary
of $J(p/q)$ to two sparse finite-group calculations.

The resulting classification has a particularly simple shape. For
mixed parity, logarithmic and rational-dilogarithmic reducibility
coincide. For odd numerator and denominator, the only nontrivial
rational-dilogarithmic cases are the ratios between $1$, $3$ and $5$.
The condition modulo $2e$, not merely modulo $e$, is essential.

A Vieta descent then parametrizes all logarithmic ratios using the
recurrences $w_{j+1}=c w_j\pm w_{j-1}$ with $c$ even. This gives both
an exact enumerator and a sparse height law. Finally, a separate analytic
argument yields a uniform formula for neighboring ratios and evaluates
the exceptional value $J(3/5)$.
These analytic identities carry their full logarithmic and $\pi^2$
terms; they are not merely boundary certificates.

\subsection{Scope of the audit and claims of novelty}

The reference snapshot is \snapshot. The main source file is
\nolinkurl{chapters/09-herglotz.tex}, whose Git blob is
\texttt{b372a62797ae9161ddcfd9f4edd036ecfd87df36}. The dependency and
open-question checks also included
\nolinkurl{chapters/10-discovery.tex} and the existing
\nolinkurl{herglotz-cyclotomic-obstructions} report.

Statements described as new are continuations relative to this reviewed
snapshot, not claims of exhaustive bibliographic priority. Classical
functional equations, the existing all-conductor symbol theorem, and
rationalized algebraic $K$-theory facts are attributed as inputs. Ordinary
proofs are supplied for the new results. The accompanying finite tests
are implementation checks, not substitutes for these proofs.

Section~\ref{sec:audit} gives one precise correction to a statement in
the cited preprint. It does not purport to audit every assertion of that
paper or of the complete ProveIt manuscript. In particular, it is not a
claim that the manuscript's independently derived boundary formula is
incorrect.

\section{The exact symbol calculus and the established kernel theorem}
\label{sec:symbols}

\subsection{Rationalized exterior symbols}

For a number field $K$, put
\[
 \A(K)=K^\times\otimes_\Z\Q,
 \qquad \cl{uv}=\cl u+\cl v.
\]
All exterior powers in this article are over $\Q$. Roots of unity vanish
in $\A(K)$. Embeddings of fields induce injective maps on these vector
spaces and on their exterior powers, so we identify old-field symbols
with their images in larger fields.

\begin{definition}[Formal reductions]\label{def:formal}
Let $\Pcal(K)$ be the rationalized pre-Bloch group: the vector space on
symbols $[z]$, modulo five-term relations, with $[0]=[1]=0$. Write
\[
 \partial[z]=\cl z\wedge\cl{1-z}\quad(z\ne0,1),
 \qquad \B(K)=\ker\partial.
\]
For a Galois-invariant element $\eta\in\Pcal(K)$, a
\emph{formal logarithmic reduction} means $\eta=0$. A
\emph{formal rational-dilogarithmic reduction} means that $\eta$ belongs
to the image of $\Pcal(\Q)$.
\end{definition}

The terminology suppresses products of logarithms of algebraic numbers
and rational multiples of $\pi^2$ in evaluation. Compatible complex
logarithms are allowed in this definition. The worked analytic identities
below use only real logarithms of positive algebraic numbers and
$\Li_2$ on $(0,1)$.

For finite Galois number fields, the standard rationalized inputs are
\begin{equation}\label{eq:K-inputs}
 \B(K)^{\Gal(K/\Q)}=\B(\Q)=0,
 \qquad \partial\Pcal(\Q)=\Lambda^2\A(\Q).
\end{equation}
These are Bloch-group descent, the torsion of $\B(\Q)$ and the torsion
of $K_2(\Q)$; the last assertion uses the identification of the
Steinberg quotient with $K_2$ for a field. They are the same inputs used
in \cite[\S7.2]{RZ}. For the Bloch-group and $K$-theoretic framework,
see \cite{ZagierDilog,CGZ}. Consequently, for invariant $\eta$,
\begin{align}
 \eta=0&\ \Longleftrightarrow\ \partial\eta=0,
 \label{eq:zero-detection}\\
 \eta\in\operatorname{im}\Pcal(\Q)
 &\ \Longleftrightarrow\
 \partial\eta\in\Lambda^2\A(\Q).
 \label{eq:rational-detection}
\end{align}
For the reverse implication in the second line, choose a rational
pre-Bloch element with the same boundary and apply the first line to
the invariant difference. This makes clear where the established
$K$-theory input enters; it is not replaced by a numerical observation.

\subsection{Cyclotomic symbols and the direct boundary formula}

For $m\ge3$ and $(a,m)=1$, define
\begin{equation}\label{eq:beta}
 \beta_m(a)=\sum_{k=1}^{m-1}
 \cl{1-\zeta_m^{ak}}\wedge\cl{1-\zeta_m^k},
 \qquad \zeta_m=e^{2\pi i/m}.
\end{equation}
Use $\beta_1=\beta_2=0$. The excluded term $k=0$ is important: there is
no multiplicative-group element $\cl0$. Reindexing gives
\begin{equation}\label{eq:symmetries}
 \beta_m(-a)=\beta_m(a),\qquad
 \beta_m(a\inv)=-\beta_m(a).
\end{equation}
Every $\beta_m(a)$ is Galois invariant. The inverse is taken modulo $m$.

For coprime positive $p,q$, the finite formula of \cite[Theorem 3]{RZ}
defines an invariant element $\xi_{p/q}\in\Pcal(\Q(\zeta_{pq}))$
representing the dilogarithmic part of $F(p/q)-F(1)$. Explicitly it is
\begin{equation}\label{eq:xi-finite}
 \xi_{p/q}=\sum_{\alpha^p=1}
 \left(\sum_{\substack{\beta^q=1\\\beta\ne1}}h(\alpha,\beta)
       -\sum_{\substack{\beta^p=1\\\beta\ne1}}h(\alpha,\beta)\right),
 \quad
 h(\alpha,\beta)=
 \left[\frac{\beta}{\beta-1}\right]
 -\left[\frac{\alpha-\beta}{1-\beta}\right].
\end{equation}

\begin{proposition}[Established boundary formula]\label{prop:boundary}
For coprime $p,q>0$,
\begin{equation}\label{eq:boundary-F}
 \partial\xi_{p/q}
 =\beta_q(p)-\beta_p(q)-\cl p\wedge\cl q.
\end{equation}
\end{proposition}
\begin{proof}
This is the boundary used in both \cite{RZ} and \cite{Repo}; we recall its
direct verification to avoid relying on the assertion audited in
Section~\ref{sec:audit}. Terms with $\alpha=1$ or $\alpha=\beta$
have zero boundary. For the others set
$A=\cl{1-\alpha}$, $B=\cl{1-\beta}$, $C=\cl{\alpha-\beta}$.
Since roots of unity vanish, direct expansion gives
\[
 \partial h(\alpha,\beta)=-C\wedge A+C\wedge B+B\wedge A.
\]
The subtracted double sum in \eqref{eq:xi-finite} has zero boundary by
antisymmetry under $\alpha\leftrightarrow\beta$. In the first double
sum, the product identities
\[
 \prod_{\substack{\beta^q=1\\\beta\ne1}}(\alpha-\beta)
 =\frac{\alpha^q-1}{\alpha-1},\qquad
 \prod_{\substack{\alpha^p=1\\\alpha\ne1}}(\beta-\alpha)
 =\frac{\beta^p-1}{\beta-1}
\]
make the first two wedge contributions $-\beta_p(q)$ and $\beta_q(p)$.
The last contribution is $\cl q\wedge\cl p$, since the products of
$1-\alpha$ and $1-\beta$ over the nontrivial roots are $p$ and $q$.
Coprimality ensures that no additional common root was omitted.
\end{proof}

\subsection{A faithful model on the symbol span}

Put $G_m=(\Z/m\Z)^\times/\{\pm1\}$ for $m\ge3$. The following theorem
is an existing result of the manuscript \cite{Repo}; its proof is
included because the new layer argument uses its Gram matrix, not just
its conclusion.

\begin{theorem}[Existing all-conductor kernel theorem]\label{thm:kernel}
For $c_a\in\Q$,
\begin{equation}\label{eq:kernel}
 \sum_{a\in G_m}c_a\beta_m(a)=0
 \quad\Longleftrightarrow\quad c_a=c_{a\inv}\quad(a\in G_m).
\end{equation}
Thus the map on the symbol span
\begin{equation}\label{eq:group-model}
 \beta_m(a)\longmapsto \ee_a-\ee_{a\inv}\in\Q[G_m]
\end{equation}
is faithful, and
\begin{equation}\label{eq:beta-zero}
 \beta_m(a)=0\quad\Longleftrightarrow\quad
 a^2\equiv\pm1\pmod m.
\end{equation}
\end{theorem}
\begin{proof}
For $1\le k<m$ define vectors and a matrix
\[
 v_k=(\log|1-\zeta_m^{bk}|)_{b\in G_m},
 \qquad A_m=\sum_{k=1}^{m-1}v_kv_k^{\mathsf T}.
\]
Let $(P_av)_b=v_{ba}$. Then $P_av_k=v_{ak}$, so $A_m$ commutes with
all $P_a$.

The vectors $v_k$ span the regular representation. Indeed their span
is invariant, and it has a nonzero projection to each character. For a
nontrivial even character of primitive conductor $f\mid m$, use
$k=m/f$. Its Fourier coefficient is a positive integer multiple of
\[
 \sum_{b\in G_f}\chi^*(b)\log|1-\zeta_f^b|
 =-\tfrac12\tau(\chi^*)L(1,\overline{\chi^*})\ne0.
\]
The equality follows by Abel-regularizing
$\log|1-r e^{i\theta}|=-\sum_{n\ge1}r^n\cos(n\theta)/n$ and applying
the primitive Gauss-sum identity. Nonvanishing uses the nonzero primitive
Gauss sum and Dirichlet's theorem $L(1,\chi)\ne0$ for nonprincipal
characters \cite[25.15.9]{DLMF}. The trivial character occurs using
$k=m/\ell$ for a prime divisor $\ell$ of $m$: the logarithmic product
is $\log\ell$; for $\ell=2$ the vector is the nonzero constant vector
$\log2$. All nontrivial roots, not just primitive $m$th roots, are used.
Therefore $A_m$ is positive definite.

Apply the exterior square of the logarithm map
$\A(\Q(\zeta_m))\to\R^{G_m}$, and identify $v\wedge w$ with
$vw^{\mathsf T}-wv^{\mathsf T}$. Its value on a symbol is
\begin{equation}\label{eq:skew-image}
 \ell_\wedge(\beta_m(a))
 =P_aA_m-A_mP_a^{\mathsf T}
 =A_m(P_a-P_{a\inv}).
\end{equation}
A zero combination of symbols therefore gives
$\sum_a(c_a-c_{a\inv})P_a=0$ after multiplication by $A_m^{-1}$.
The regular permutation matrices are linearly independent, so the
coefficients are symmetric. Conversely symmetric coefficients cancel
by \eqref{eq:symmetries}; self-inverse classes contribute zero over
$\Q$. This proves all assertions.
\end{proof}

\begin{remark}
The exterior logarithm map is not asserted to be injective on all of
$\Lambda^2\A(K)$. Its explicitly computed value suffices on the
particular span in Theorem~\ref{thm:kernel}. Positivity above is proved
by character nonvanishing, not by a floating-point determinant.
\end{remark}

\subsection{Separating coprime conductors}

\begin{lemma}[Normalized norm separation]\label{lem:norm}
Let $(r,s)=1$, with $U$ a rational combination of the $\beta_r(a)$ and
$V$ a rational combination of the $\beta_s(b)$, viewed in a common
cyclotomic field. Then
\[
 U-V\in\Lambda^2\A(\Q)\quad\Longleftrightarrow\quad U=V=0.
\]
Moreover, normalized norm to $\Q$ kills each such symbol combination.
\end{lemma}
\begin{proof}
The linear normalized field norm to $E$ is
$N_E\cl u=\cl{\Norm_{K/E}u}/[K:E]$. It fixes old-field elements;
apply its exterior square to wedges. Coprime cyclotomic fields are
linearly disjoint over $\Q$. The norm of either factor of a term of
$\beta_s(b)$ to $\Q(\zeta_r)$ is the same rational vector, since the
factors are Galois conjugates in $\Q(\zeta_s)$. Thus that wedge is
killed, while every old $r$-conductor wedge is fixed. Reversing the
roles gives the other projection. If $U-V=W$ is rational, the two
projections give $U=W$ and $-V=W$; projecting either again to $\Q$
gives $W=0$. The cases of conductor $1$ or $2$ follow from the
zero-symbol conventions. The same conjugate-factor argument proves
the last assertion.
\end{proof}

\section{Conductor doubling and rigidity of the new dyadic layer}
\label{sec:dyadic}

This section contains the new structural input. Doubling an odd conductor
does not enlarge its cyclotomic field, whereas doubling an even conductor
does. The two cases must be treated differently.

\subsection{Odd-conductor doubling}

\begin{theorem}[Exact conductor-doubling identity]\label{thm:doubling}
Let $m$ be odd and $(a,2m)=1$. In
$\Lambda^2\A(\Q(\zeta_m))$,
\begin{equation}\label{eq:doubling}
 \beta_{2m}(a)=3\beta_m(a)-\beta_m(2a)-\beta_m(a/2),
\end{equation}
where $a/2$ denotes multiplication by the inverse of $2$ modulo $m$.
The formula also holds for $m=1$ under the zero-symbol convention.
\end{theorem}
\begin{proof}
For $m>1$ split $\mu_{2m}$ into $\mu_m$ and $-\mu_m$. The first
part of the sum defining $\beta_{2m}(a)$ is $\beta_m(a)$. Set
$u_k=\cl{1-\zeta_m^k}$ for $k\not\equiv0\pmod m$. For these $k$,
\[
 \cl{1+\zeta_m^k}=u_{2k}-u_k.
\]
Since $a$ is odd, the second part, apart from the zero wedge at the root
$-1$, is
\[
 \sum_{k=1}^{m-1}(u_{2ak}-u_{ak})\wedge(u_{2k}-u_k).
\]
The two diagonal terms each give $\beta_m(a)$ by reindexing. The cross
terms give $-\beta_m(2a)$ and $-\beta_m(a/2)$. Including the first
part proves \eqref{eq:doubling}. No logarithm branch is used in this
multiplicative-group calculation.
\end{proof}

\subsection{Even-conductor descent is rigid}

\begin{theorem}[New dyadic layer rigidity]\label{thm:layer}
Let $m$ be even and $(a,2m)=1$. Then
\begin{equation}\label{eq:layer}
 \beta_{2m}(a)\in\Lambda^2\A(\Q(\zeta_m))
 \quad\Longleftrightarrow\quad
 \beta_{2m}(a)=0
 \quad\Longleftrightarrow\quad
 a^2\equiv\pm1\pmod{2m}.
\end{equation}
\end{theorem}
\begin{proof}
For $m=2$, all the symbols at conductor $4$ vanish, so the assertion
is immediate. Suppose $m\ge4$. The natural map
$G_{2m}\to G_m$ has kernel $\{1,h\}$, where $h$ is the class of
$1+m$. This is a nontrivial element of order $2$.

In the real regular representation $\R^{G_{2m}}$, logarithm vectors
coming from $\Q(\zeta_m)$ lie in the $+1$ eigenspace of $P_h$.
Consequently, a skew matrix arising from an exterior product of such
vectors annihilates the $-1$ eigenspace. If $\beta_{2m}(a)$ descends,
\eqref{eq:skew-image} therefore gives
\[
 A_{2m}(P_a-P_{a\inv})(I-P_h)=0.
\]
The Gram matrix is invertible. The faithful regular representation of
the group algebra then gives
\begin{equation}\label{eq:four-terms}
 (\ee_a-\ee_{a\inv})(\ee_1-\ee_h)=0.
\end{equation}
If $a\ne a\inv$ in $G_{2m}$, the coefficient of $\ee_a$ in the left
side is positive: the negative term $-\ee_{a\inv}$ is distinct,
$-\ee_{ah}$ is distinct because $h\ne1$, and any coincidence with
$+\ee_{a\inv h}$ only increases that coefficient. Hence
\eqref{eq:four-terms} is impossible. Thus $a=a\inv$ in $G_{2m}$,
which is exactly the last congruence in \eqref{eq:layer}.
Theorem~\ref{thm:kernel} now makes the symbol zero. The converse
implication to descent is immediate.
\end{proof}

The theorem is stronger than merely testing the normalized norm of the
new symbol. It detects a nonzero new-character component, so no
combination of old-field wedges can cancel it. This is the obstruction
that forces a modulus $2e$ in the final classification.

\subsection{Two finite-group difference tests}

For an odd conductor $m$, division by $2$ in symbol arguments is
well-defined. The following two tests will isolate the denominator and
numerator contributions of $J$.

\begin{lemma}[First difference]\label{lem:first}
For odd $m\ge1$ and $(a,m)=1$,
\begin{equation}\label{eq:first}
 \beta_m(a)-\beta_m(a/2)=0
 \quad\Longleftrightarrow\quad m\in\{1,3,5\}.
\end{equation}
\end{lemma}
\begin{proof}
The case $m=1$ is immediate. For $m\ge3$, put $x=\overline a$ and
$b=\overline2$ in $G_m$, and let $f_x=\ee_x-\ee_{x\inv}$.
By Theorem~\ref{thm:kernel}, the required equality is
$f_x=f_{x/b}$. If this vector is nonzero, its unique positive basis
coefficient identifies its index, so $x=x/b$ and $b=1$. If both
vectors are zero, $x^2=(x/b)^2=1$, so $b^2=1$. In either case
$b^2=1$, or $4\equiv\pm1\pmod m$, forcing odd $m$ to be $1$, $3$
or $5$. Conversely the groups at $3$ and $5$ have exponent at most
$2$, and all their symbols vanish.
\end{proof}

\begin{lemma}[Second difference]\label{lem:second}
For odd $m\ge1$ and $(a,m)=1$,
\begin{equation}\label{eq:second}
 \beta_m(2a)-2\beta_m(a)+\beta_m(a/2)=0
 \quad\Longleftrightarrow\quad a^2\equiv\pm1\pmod m.
\end{equation}
\end{lemma}
\begin{proof}
Again $m=1$ is automatic. In the group algebra of $G_m$, the image of
the left side is $(T_b+T_{b\inv}-2I)f_x$, where $T_b$ is translation
by $b=\overline2$. With the orthonormal basis $\ee_g$,
\[
 \big\langle f,(2I-T_b-T_{b\inv})f\big\rangle
 =\|f-T_bf\|^2.
\]
Vanishing of the second difference therefore forces $f_x=T_bf_x$.
If $f_x\ne0$, the unique positive coefficient of $f_x$ forces $b=1$.
For an odd conductor this means $m=1$ or $3$, at which $f_x=0$
already, a contradiction. Thus $f_x=0$, and the kernel theorem gives
the asserted congruence. Conversely $f_x=0$ makes its second difference
zero. The two positive and two negative translated terms agree with
the symbol expression even though a single translation of $f_x$ need
not itself have the form $f_y$.
\end{proof}

\section{Complete formal classification at rational arguments}
\label{sec:classification}

For positive rational $x$, reduce every subscript to lowest terms and set
\begin{equation}\label{eq:eta}
 \eta_x=\xi_{2x}-2\xi_x+\xi_{x/2}.
\end{equation}
By \eqref{eq:J-F}, this is the dilogarithmic element relevant to $J(x)$.
The extra rational multiple of $\pi^2$ has no boundary.

\begin{lemma}[Analytic and formal reciprocity]\label{lem:reciprocity}
For $x>0$,
\begin{equation}\label{eq:J-reciprocity}
 J(x)+J(1/x)=\log^22.
\end{equation}
For positive rational $x$, one also has $\eta_{1/x}=-\eta_x$ in a
common pre-Bloch group. Hence each type of formal reduction is invariant
under $x\mapsto1/x$.
\end{lemma}
\begin{proof}
Integration by parts gives
\[
 J(x)=\log^22-\int_0^1
       \log(1+t)\frac{x t^{x-1}}{1+t^x}\dd t.
\]
The substitution $u=t^x$ identifies the integral with $J(1/x)$. This
known reciprocity also follows from the functional equations discussed
in \cite{CK}.

For the formal statement, \eqref{eq:boundary-F} gives
$\partial\xi_{1/r}=-\partial\xi_r$. Their invariant sum is zero by
\eqref{eq:zero-detection}. Applying this to the three terms of
\eqref{eq:eta} proves the assertion. Thus the formal assertion is not
inferred solely from an equality of numerical periods.
\end{proof}

\begin{theorem}[Complete all-rational classification]\label{thm:classification}
Let $p,q$ be coprime positive integers.

\smallskip\noindent
\textbf{Mixed parity.} Let $e$ denote the even entry and $o$ the odd entry
of $(p,q)$. Then $J(p/q)$ has a formal rational-dilogarithmic reduction
if and only if it has a formal logarithmic reduction, and these are
equivalent to
\begin{equation}\label{eq:mixed-class}
 \boxed{\quad e^2\equiv\pm1\pmod o,
 \qquad o^2\equiv1\pmod{2e}.\quad}
\end{equation}
Congruences modulo $1$ are interpreted as automatic.

\smallskip\noindent
\textbf{Both entries odd.} There is a formal rational-dilogarithmic
reduction exactly when
\begin{equation}\label{eq:odd-class}
 p,q\in\{1,3,5\}.
\end{equation}
There is a formal logarithmic reduction exactly when $p=q=1$.

\smallskip\noindent
Consequently, the rational-dilogarithmic but not logarithmic arguments
are precisely
\begin{equation}\label{eq:six-extras}
 \boxed{\quad\frac13,\ 3,\ \frac15,\ 5,\ \frac35,\ \frac53.\quad}
\end{equation}
\end{theorem}

\begin{proof}
For mixed parity, Lemma~\ref{lem:reciprocity} lets us take $x=e/o$.
Expanding \eqref{eq:boundary-F} gives
\begin{equation}\label{eq:mixed-boundary}
 \partial\eta_{e/o}=D_o(e)-N_e(o),
\end{equation}
where
\begin{align}
 D_o(e)&=\beta_o(2e)-2\beta_o(e)+\beta_o(e/2),\label{eq:D}\\
 N_e(o)&=\beta_{2e}(o)-2\beta_e(o)+\beta_{e/2}(o).\label{eq:N}
\end{align}
The rational wedges cancel because
$-\cl{2e}+2\cl e-\cl{e/2}=0$. The conductors $o$ and $2e$ are
coprime. By Lemma~\ref{lem:norm}, the boundary is rational if and only
if both $D_o(e)$ and $N_e(o)$ vanish. It is then zero, proving that
rational-dilogarithmic and logarithmic reducibility coincide.

Lemma~\ref{lem:second} gives
$D_o(e)=0\Longleftrightarrow e^2\equiv\pm1\pmod o$.
If $N_e(o)=0$, its top-conductor term belongs to
$\Lambda^2\A(\Q(\zeta_e))$. Theorem~\ref{thm:layer} forces
$\beta_{2e}(o)=0$, or $o^2\equiv\pm1\pmod{2e}$.
Since $4\mid2e$ and $o$ is odd, the negative sign is impossible.
Conversely $o^2\equiv1\pmod{2e}$ kills all three symbols of
$N_e(o)$, including the conductor-$1$ or conductor-$2$ degeneracies.
This proves \eqref{eq:mixed-class}.

Now suppose both $p$ and $q$ are odd. Applying
Theorem~\ref{thm:doubling} at both conductors to the boundary expansion
yields
\begin{equation}\label{eq:odd-boundary}
 \begin{split}
 \partial\eta_{p/q}={}&
  \bigl(\beta_q(p)-\beta_q(p/2)\bigr)
 -\bigl(\beta_p(q)-\beta_p(q/2)\bigr)\\
 &+\cl2\wedge\cl{p/q}.
 \end{split}
\end{equation}
For example, the rational contribution before simplification is
$-\cl{2p}\wedge\cl q+2\cl p\wedge\cl q
 -\cl p\wedge\cl{2q}$, which is the last term displayed.
Normalized norm separation and Lemma~\ref{lem:first} show that this
boundary is rational exactly when $p,q\in\{1,3,5\}$.

For logarithmic reduction, normalized norm to $\Q$ first requires
$\cl2\wedge\cl{p/q}=0$. If the reduced odd ratio $p/q$ is not $1$,
the prime-exponent vector of $p/q$ is nonzero and has zero coordinate
at the prime $2$. It is therefore independent of $\cl2$ in
$\A(\Q)$, making this wedge nonzero. If $p=q=1$, all terms vanish
and $J(1)=\tfrac12\log^22$. This proves the odd case and the list
\eqref{eq:six-extras}.
\end{proof}

\subsection{Examples and the missing factor of two}

The mixed pairs $(e,o)=(4,15),(12,5),(12,29)$ satisfy
\eqref{eq:mixed-class}. Thus $J(4/15)$, $J(5/12)$ and $J(12/29)$ have
formal logarithmic reductions. Explicit short logarithmic expressions
for arbitrary pairs of this kind are a separate constructive task;
the neighboring-ratio family will be evaluated below.

By contrast $(e,o)=(8,3)$ satisfies
$8^2\equiv1\pmod3$ and $3^2\equiv1\pmod8$, but fails
$3^2\equiv1\pmod{16}$. The new-layer obstruction can be seen exactly
in $G_{16}$: with $h=9$, the four-term expression
\eqref{eq:four-terms} at $a=3$ is
\[
 (\ee_3-\ee_5)(\ee_1-\ee_7)=2\ee_3-2\ee_5\ne0,
\]
where residues have been identified up to sign. Hence $J(8/3)$ has no
formal rational-dilogarithmic reduction. Testing only modulo $e$ would
produce a false positive.

For $x=2/q$ the theorem recovers exactly the existing classification.
Odd $q$ gives $q=3,5$; a denominator divisible by $4$ reduces to a ratio
with even entry and odd entry $1$; and a denominator $q=2m$ with odd
$m$ gives the odd ratio $1/m$, allowing precisely $m=1,3,5$ for rational
dilogarithms. Thus the positive cases are $q=2,3,5,6,10$ or $4\mid q$;
the cases $6,10$ are not logarithmic in the formal calculus.

\begin{remark}[What has not been proved]\label{rem:periods}
Theorem~\ref{thm:classification} is a necessary and sufficient condition
for a five-term certificate in Definition~\ref{def:formal}. For a ratio
that fails the test, it does not prove that the real number $J(p/q)$ is
linearly independent of all logarithmic products or rational dilogarithm
values. Such a claim would require information about numerical period
evaluation beyond the symbol theorem. In particular, the numerical
transcendence problems in the manuscript remain separate.
\end{remark}

\section{Even quadratic recurrences and height asymptotics}
\label{sec:recurrences}

The modular criterion admits a complete elementary parametrization.
It also gives a way to enumerate the exceptional rational set without
factoring either entry.

\begin{theorem}[Recurrence parametrization]\label{thm:recurrences}
Apart from $J(1)$, all unordered reduced logarithmic pairs $1\le a<b$
are exactly the adjacent pairs in the following families, for every even
integer $c\ge2$:
\begin{align}
 u_0=1,\quad u_1=c,\quad u_{j+1}&=c u_j-u_{j-1},\label{eq:minus-rec}\\
 v_0=1,\quad v_1=c,\quad v_{j+1}&=c v_j+v_{j-1}.\label{eq:plus-rec}
\end{align}
Both $a/b$ and $b/a$ are logarithmic cases. The seed $(1,c)$ is common
to the two families. Beyond the seeds there is no overlap between the
families or between different values of $c$.
\end{theorem}
\begin{proof}
Let $e$ be even and $o$ odd in a pair satisfying
\eqref{eq:mixed-class}. First suppose $e^2\equiv1\pmod o$.
Then
\[
 c=\frac{e^2+o^2-1}{eo}
\]
is an integer. It is even: after division of the numerator by $e$, both
$e$ and $(o^2-1)/e$ are even, and the remaining divisor $o$ is odd.
Also $c\ge2$, since $(e-o)^2\ge1$. Writing the sorted pair as $a<b$
gives
\begin{equation}\label{eq:minus-invariant}
 a^2+b^2-cab=1.
\end{equation}
If $a>1$, the new entry
$d=(a^2-1)/b=ca-b$ is a positive integer strictly smaller than $a$.
The pair $(d,a)$ has the same invariant, is coprime, and has opposite
parity because $c$ is even and $d$ has the parity of $b$. The invariant
itself gives both congruences in \eqref{eq:mixed-class}. Repetition
reaches a pair $(1,c)$. Reversing the descent is
\eqref{eq:minus-rec}.

Next suppose $e^2\equiv-1\pmod o$. The signed integer
\[
 H=\frac{o^2-e^2-1}{eo}
\]
is even by the same divisibility argument, and is nonzero: otherwise
$(o-e)(o+e)=1$, impossible for $e>0$. Put $c=|H|\ge2$. In sorted
coordinates,
\begin{equation}\label{eq:plus-invariant}
 b^2-cab-a^2=\varepsilon,
 \qquad \varepsilon=\begin{cases}1,&b\text{ odd},\\-1,&b\text{ even}.
 \end{cases}
\end{equation}
For $a>1$ set
$d=b-ca=(a^2+\varepsilon)/b$. This is a positive integer less than
$a$; for $\varepsilon=1$, use $a^2+1<a(a+1)\le ab$. The pair
$(d,a)$ is coprime and has opposite parity. Its invariant has sign
$-\varepsilon$, exactly as required when the parity of the larger entry
changes. Thus the descent continues to $(1,c)$, and its reverse is
\eqref{eq:plus-rec}.

Conversely the recurrences preserve \eqref{eq:minus-invariant} or
\eqref{eq:plus-invariant}, respectively, and preserve coprimality and
alternating parity. In the minus family the invariant gives
$e^2\equiv1\pmod o$ and
$o^2-1=e(co-e)$, which is divisible by $2e$. In the plus family it
gives $e^2\equiv-1\pmod o$ and
$o^2-1=e(e\pm co)$, again divisible by $2e$. Thus every generated
pair satisfies the criterion.

If the odd entry exceeds $1$, the signs $+1$ and $-1$ modulo that
entry cannot both hold. The displayed invariants determine $c$
uniquely within each case, and the strictly decreasing descent
identifies a unique depth. The only possible shared pairs are the seeds,
whose odd entry is $1$.
\end{proof}

\begin{center}
\begin{tabular}{lll}
\toprule
Family & Coefficient & Initial terms\\
\midrule
Minus & $2$ & $1,2,3,4,5,6,\ldots$\\
Minus & $4$ & $1,4,15,56,209,\ldots$\\
Minus & $6$ & $1,6,35,204,1189,\ldots$\\
Plus & $2$ & $1,2,5,12,29,70,\ldots$\\
Plus & $4$ & $1,4,17,72,305,\ldots$\\
Plus & $6$ & $1,6,37,228,1405,\ldots$\\
\bottomrule
\end{tabular}
\end{center}

The coefficient-$2$ minus family explains why every neighboring ratio
is logarithmic. All other families grow geometrically. In particular,
a descent algorithm can shortcut consecutive integers and otherwise
uses $O(\log\max(a,b))$ steps. This statement counts arithmetic descent
steps; it is not a bound on the bit complexity of an arbitrary expanded
analytic certificate.

\subsection{A complete hierarchy of height terms}

Let $N(B)$ count unordered coprime pairs $1\le a<b\le B$ satisfying the
logarithmic criterion; $J(1)$ is not counted. Here $B$ is a positive
integer.

\begin{theorem}[Height expansion to every fixed order]\label{thm:counting}
For every fixed integer $R\ge2$, as $B\to\infty$,
\begin{equation}\label{eq:count-expansion}
 N(B)=\frac32 B+\sum_{r=2}^{R}B^{1/r}
       +O_R\bigl(B^{1/(R+1)}\log B\bigr).
\end{equation}
In particular,
\begin{equation}\label{eq:count-two}
 N(B)=\frac32 B+\sqrt B+O(B^{1/3}\log B).
\end{equation}
The number of ordered positive rational logarithmic arguments of height
at most $B$ is $2N(B)+1$. For $B\ge5$, the number admitting a formal
rational-dilogarithmic reduction is $2N(B)+7$.
\end{theorem}
\begin{proof}
Count each seed $(1,c)$ once: there are $\lfloor B/2\rfloor$.
The minus family at coefficient $2$ contributes all consecutive pairs;
after removing its seed it contributes $B-2$ for $B\ge2$. Their sum
is $\tfrac32B+O(1)$.

For the other pairs write $w_r^{\pm}(c)$ for the term at index $r$ in
the two recurrences. The pair with this height has indices $(r-1,r)$.
At fixed $r\ge2$, induction on the recurrence gives
\[
 w_r^{\pm}(c)=c^r\pm(r-1)c^{r-2}+O_r(c^{r-4}),
\]
where the remainder can simply be read as the lower-degree polynomial
terms; when $r$ is small it is zero as appropriate. Equivalently,
$w_r^{\pm}(c)=c^r+O_r(c^{r-2})$. The number of even coefficients
satisfying $w_r^{\pm}(c)\le B$ is therefore
$\tfrac12 B^{1/r}+O_r(1)$: the transition can differ from
$c=B^{1/r}$ by only $O_r(B^{-1/r})$, with finitely many small
coefficients absorbed in the constant. Removing coefficient $2$ from
the minus family changes this by $O(1)$. Both signs together contribute
$B^{1/r}+O_r(1)$. Summing over $2\le r\le R$ gives the displayed
fractional powers.

For completeness, the remaining depths are uniformly bounded. In the
minus family with $c\ge4$, induction on the successive ratios gives
$w_j/w_{j-1}>c-1$, so $w_j\ge(c-1)^j$. In the plus family with
$c\ge2$, $w_j\ge c^j$. A term with index at least $R+1$ and height
at most $B$ therefore requires
$c\le 1+B^{1/(R+1)}$. There are
$O(B^{1/(R+1)})$ possible coefficients, and each contributes only
$O(\log B)$ depths, because its growth factor is at least $2$.
This proves the remainder estimate. Disjointness from
Theorem~\ref{thm:recurrences} prevents overcounting.

Finally each unordered pair yields two reciprocal arguments, while
$x=1$ adds one. The six distinct extra arguments in
\eqref{eq:six-extras} all have height at most $5$.
\end{proof}

The theorem implies a sparse set among all reduced rational arguments
of bounded height, but no transcendence statement is attached to that
sparsity. It counts exactly the formal reduction cases.

\begin{center}
\begin{tabular}{r r r}
\toprule
$B$ & $N(B)$ & Ordered logarithmic arguments\\
\midrule
$100$ & $161$ & $323$\\
$1\,000$ & $1\,543$ & $3\,087$\\
$10\,000$ & $15\,135$ & $30\,271$\\
$1\,000\,000$ & $1\,501\,162$ & $3\,002\,325$\\
$100\,000\,000$ & $150\,010\,658$ & $300\,021\,317$\\
$10\,000\,000\,000$ & $15\,000\,102\,700$ & $30\,000\,205\,401$\\
\bottomrule
\end{tabular}
\end{center}
These are exact counts from the recurrence enumerator, not floating-point
fits to \eqref{eq:count-expansion}. The script uses integer arithmetic
and does not enumerate all $O(B^2)$ pairs at the larger heights.

\section{Explicit analytic identities with full constants}
\label{sec:identities}

The formal classification establishes existence of reductions but does
not, by itself, print their logarithmic corrections. This section does
so for an infinite family and for the remaining exceptional ratio
$3/5$. All the identities in this section are ordinary analytic
equalities. No integer-relation search is used in their proofs.

\subsection{Classical normalized functional equations}

For $0<t<1$ define the Rogers dilogarithm
\begin{equation}\label{eq:Rogers}
 \Rog(t)=\Li_2(t)+\frac12\log t\log(1-t),
 \qquad \Rog(1/2)=\frac{\pi^2}{12}.
\end{equation}
We use its reflection and five-term identities in the real interval:
\begin{align}
 \Rog(t)+\Rog(1-t)&=\frac{\pi^2}{6},\label{eq:Rog-reflect}\\
 \Rog(A)+\Rog(B)&=\Rog(AB)
  +\Rog\left(\frac{A(1-B)}{1-AB}\right)
  +\Rog\left(\frac{B(1-A)}{1-AB}\right),
 \quad 0<A,B<1.\label{eq:Rog-five}
\end{align}
These are classical; their normalization and relation to the ordinary
dilogarithm are reviewed in \cite{ZagierDilog,RZ}.

Put
\begin{equation}\label{eq:modified-F}
 f(x)=F(x)-F(1)+\frac{\pi^2}{12}(x-2+x^{-1})-\frac14\log^2x.
\end{equation}
The functional equations needed below are the normalized inversion and
three-term equations and the determinant-$2$ Hecke equation of
\cite[(15), (22a)]{RZ}:
\begin{align}
 f(1/x)&=-f(x),\label{eq:f-inversion}\\
 f(x)-f(x+1)-f\left(\frac{x}{x+1}\right)
 &=\Rog(1/2)-\Rog\left(\frac{x}{x+1}\right),\label{eq:f-three}\\
 \begin{split}
 f(2x)+f(x/2)+f\left(\frac{x+1}{2}\right)
 +f\left(\frac{2x}{x+1}\right)-3f(x)
 ={}&\Rog\left(\frac{x}{x+1}\right)-\Rog(1/2)\\
 &+\tfrac12\log2\log x.
 \end{split}\label{eq:f-Hecke}
\end{align}
They hold for $x>0$ with real logarithms. Formula \eqref{eq:J-F}
becomes
\begin{equation}\label{eq:J-f}
 J(x)=f(2x)-2f(x)+f(x/2)+\frac12\log^22
       -\frac{\pi^2}{24}(x-x^{-1}).
\end{equation}

Define the finite log-sine sum
\begin{equation}\label{eq:S}
 S(m)=\sum_{k=1}^{m-1}\log^2\left(2\sin\frac{\pi k}{m}\right),
 \qquad S(1)=0.
\end{equation}
Every sine argument in this sum is positive. The integer formula of
\cite[(23)]{RZ} simplifies under \eqref{eq:modified-F} to
\begin{equation}\label{eq:f-integer}
 f(m)=\frac{(m-1)\pi^2}{24}+\frac12S(m)+\frac14\log^2m,
 \qquad m\ge1.
\end{equation}
These classical inputs are kept separate from the deductions that follow.

\subsection{Every neighboring rational argument}

\begin{theorem}[Uniform neighboring-ratio identity]\label{thm:neighbor}
For every integer $n\ge1$,
\begin{equation}\label{eq:neighbor}
 \boxed{\begin{aligned}
 J\left(\frac n{n+1}\right)={}&\frac12\log^22
 +\frac12\bigl(S(n)-S(n+1)+S(2n+2)-S(2n)\bigr)\\
 &-\frac{\pi^2(n^2-n-1)}{24n(n+1)}.
 \end{aligned}}
\end{equation}
The inverse ratios are evaluated by \eqref{eq:J-reciprocity}.
\end{theorem}
\begin{proof}
Solve \eqref{eq:f-Hecke} for the second difference in
\eqref{eq:J-f}. Set
\[
 x=\frac n{n+1},\quad A=\frac{2n+1}{2n+2},\quad
 B=\frac{2n}{2n+1},\quad C=\frac n{2n+1}.
\]
The three $f$-terms on the resulting right side are $f(x)-f(A)-f(B)$.
For any integer $m\ge1$, \eqref{eq:f-three} gives
\[
 f\left(\frac m{m+1}\right)=f(m)-f(m+1)-\Rog(1/2)
                                      +\Rog\left(\frac m{m+1}\right).
\]
The integer terms telescope, leaving
\begin{equation}\label{eq:neighbor-intermediate}
 \begin{split}
 f(2x)-2f(x)+f(x/2)={}&f(n)-f(n+1)+f(2n+2)-f(2n)\\
 &+\Rog(x)-\Rog(A)-\Rog(B)+\Rog(C)
   +\tfrac12\log2\log x.
 \end{split}
\end{equation}
In \eqref{eq:Rog-five} with these $A,B$, the three right-hand arguments
are respectively $x$, $1/2$ and $C$. Hence the Rogers combination in
\eqref{eq:neighbor-intermediate} is exactly $-\Rog(1/2)$.

Substitution of \eqref{eq:f-integer} contributes $\pi^2/24$ from the
four integer terms. Their ordinary log-square contribution is
\[
 \tfrac14\bigl(\log^2n-\log^2(n+1)+\log^2(2n+2)-\log^2(2n)\bigr)
 =-\tfrac12\log2\log x,
\]
which cancels the cross term. The remaining $\pi^2$ coefficient in
\eqref{eq:J-f} is
\[
 \frac1{24}-\frac1{12}-\frac1{24}(x-x^{-1})
 =-\frac{n^2-n-1}{24n(n+1)}.
\]
This proves \eqref{eq:neighbor}. All Rogers arguments used in the proof
are strictly between zero and one, so no branch adjustment or unspecified
constant is hidden in the cancellation.
\end{proof}

\begin{corollary}[Small explicit logarithmic evaluations]\label{cor:small-J}
Let $\phi=(1+\sqrt5)/2$. Then
\begin{align}
 J(4/3)&=\frac78\log^22-\frac12\log^2(1+\sqrt2)
                     +\frac{5\pi^2}{288},\label{eq:J43}\\
 J(4/5)&=\frac78\log^22+2\log^2\phi
              -\frac12\log^2(1+\sqrt2)-\frac{11\pi^2}{480}.
 \label{eq:J45}
\end{align}
\end{corollary}
\begin{proof}
Write $a=\log2$, $b=\log3$, $c=\log5$, $d=\log\phi$ and
$u=\log(1+\sqrt2)$. Elementary trigonometric products give
\begin{align*}
 S(3)&=\tfrac12b^2,&S(4)&=\tfrac32a^2,&
 S(5)&=\tfrac14c^2+d^2,\\
 S(6)&=a^2+\tfrac12b^2,&S(8)&=\tfrac74a^2+u^2,&
 S(10)&=a^2+\tfrac14c^2+5d^2.
\end{align*}
For instance $2\sin(\pi/10)=\phi^{-1}$ and
$2\sin(3\pi/10)=\phi$, while the two distinct primitive positive
logarithms at conductor $8$ have sum $a/2$ and difference $u$.
Substitute into \eqref{eq:neighbor} for $n=3,4$, using reciprocity for
$n=3$.
\end{proof}

At $n=1,2$, the theorem also recovers
$J(1/2)=\pi^2/48+\log^22/4$ and
$J(2/3)=-\pi^2/144+3\log^22/4$. These are consistency checks against
known special cases, not additional novelty claims.

\subsection{The exceptional value at \texorpdfstring{$3/5$}{3/5}}

\begin{theorem}[A compact exceptional dilogarithm identity]\label{thm:J35}
With $\phi=(1+\sqrt5)/2$,
\begin{equation}\label{eq:J35}
 \boxed{\begin{aligned}
 J(3/5)={}&\Li_2(1/3)-\frac12\Li_2(1/5)
  +\frac12\log^22+\frac12\log^23-\frac14\log^25\\
 &+\log^2\phi-\frac{7\pi^2}{180}.
 \end{aligned}}
\end{equation}
In particular $[1/3]-\tfrac12[1/5]$ is a rational-dilogarithm core.
\end{theorem}
\begin{proof}
We first track the normalized Herglotz terms, then remove the remaining
Rogers terms by a single five-term substitution. Set
$a=\log2$, $b=\log3$, $c=\log5$, $d=\log\phi$.
The classical $F(2/5)$ identity \cite[(24)]{RZ} is
\[
 F(2/5)-F(1)=\tfrac12\Li_2(4/5)-\tfrac15\pi^2
 +\log^2(2\sin(\pi/5))+\log^2(2\sin(2\pi/5))
 +a\log(2/5).
\]
The two log-sine squares sum to $c^2/8+d^2/2$. Converting to $f$ and
to the Rogers normalization gives
\begin{equation}\label{eq:f25}
 f(2/5)=\tfrac12\Rog(4/5)-\tfrac18\pi^2
                      +\tfrac34a^2-\tfrac38c^2+\tfrac12d^2.
\end{equation}
The integer formula gives
\begin{align*}
 f(2)&=\pi^2/24+3a^2/4,&
 f(3)&=\pi^2/12+b^2/2,\\
 f(4)&=\pi^2/8+7a^2/4,&
 f(5)&=\pi^2/6+3c^2/8+d^2/2.
\end{align*}
Therefore
\begin{equation}\label{eq:f25-collapse}
 -f(2)+f(5)+f(2/5)=d^2+\tfrac12\Rog(4/5).
\end{equation}
Inversion and the three-term relation at $3/2$ give
\[
 f(3/5)=-f(2)+f(3)+f(2/5)-\Rog(2/3)+\Rog(3/5).
\]
Also
$f(4/5)=f(4)-f(5)-\Rog(1/2)+\Rog(4/5)$ and
$f(3/4)=f(3)-f(4)-\Rog(1/2)+\Rog(3/4)$.
The Hecke relation at $x=3/5$, together with
\eqref{eq:J-f} and \eqref{eq:f25-collapse}, now becomes
\begin{equation}\label{eq:J35-Rog}
 J(3/5)=d^2+T+\tfrac12a(b-c)+\tfrac12a^2+\frac{2\pi^2}{45},
\end{equation}
where
\[
 T=\Rog(1/2)-\Rog(2/3)+\Rog(3/5)-\tfrac12\Rog(4/5)
                         -\Rog(3/4)+\Rog(3/8).
\]
The substitution $A=1/2$, $B=3/4$ into \eqref{eq:Rog-five} gives
\[
 \Rog(1/2)+\Rog(3/4)=\Rog(3/8)+\Rog(1/5)+\Rog(3/5).
\]
Combining this with reflection at $1/3$ and $1/5$ yields
\[
 T=\Rog(1/3)-\tfrac12\Rog(1/5)-\frac{\pi^2}{12}.
\]
Finally,
\[
 \Rog(1/3)-\tfrac12\Rog(1/5)
 =\Li_2(1/3)-\tfrac12\Li_2(1/5)
  -\tfrac12ab+\tfrac12b^2+\tfrac12ac-\tfrac14c^2.
\]
The cross terms cancel in \eqref{eq:J35-Rog};
$2/45-1/12=-7/180$ gives the stated constant.

To check the core directly, take its boundary:
\[
 \partial\bigl([1/3]-\tfrac12[1/5]\bigr)
 =\cl2\wedge\cl3-\cl2\wedge\cl5
 =\cl2\wedge\cl{3/5},
\]
which agrees with \eqref{eq:odd-boundary}. The analytic argument above,
not this last boundary check alone, supplies the complete constants.
\end{proof}

\subsection{All six non-logarithmic formal exceptions are explicit}

The known identities for the other two reciprocal pairs are
\begin{align}
 J(1/3)&=\frac{\pi^2}{9}+\frac12\log^22-\frac12\log^23-\Li_2(1/3),
 \label{eq:J13}\\
 J(1/5)&=\frac{7\pi^2}{60}+\frac12\log^22-\frac14\log^25
                  -\log^2\phi-\frac12\Li_2(1/5).
 \label{eq:J15}
\end{align}
These are already present in the manuscript's $J(2/q)$ treatment
\cite{Repo}. For a direct derivation, use \eqref{eq:J-F} at $1/3$ and
$1/5$, the reciprocal-integer formula
\[
 F(1/m)-F(1)=-\frac{(m-1)(3m-2)}{24m}\pi^2-\frac12S(m),
\]
the three-term consequence
\[
 F(2/3)-F(1)=-\frac{5\pi^2}{36}+\Li_2(2/3)
                  +\frac12\log^2(3/2)-\frac12S(3)+\frac12S(2),
\]
and the displayed $F(2/5)$ formula. The elementary values of $S$ and
Euler reflection at $1/3$ and $1/5$ give
\eqref{eq:J13}--\eqref{eq:J15} without a branch change.

Equations \eqref{eq:J35}, \eqref{eq:J13}, \eqref{eq:J15} and
reciprocity evaluate every argument in \eqref{eq:six-extras}. Their
completeness as a list of formal exceptions is the new classification;
the first two classical evaluations are not newly asserted discoveries.

\begin{corollary}[Cancellation among the three exceptional ratios]
\label{cor:exceptional-relation}
\begin{equation}\label{eq:exceptional-relation}
 \boxed{\quad J(1/3)-J(1/5)+J(3/5)
      =\frac12\log^22+2\log^2\phi-\frac{2\pi^2}{45}.\quad}
\end{equation}
\end{corollary}
\begin{proof}
Substitute the three evaluations. Both rational dilogarithms, the
$\log^23$ term and the $\log^25$ term cancel. The rational
$\pi^2$ coefficient is $1/9-7/60-7/180=-2/45$.
\end{proof}

\section{An exact correction to a finite-field assertion}
\label{sec:audit}

The reviewed arXiv preprint of Radchenko--Zagier, version
\texttt{2012.15805v1}, Section~7.2, printed page~16, includes the
finite-field assertion
\begin{equation}\label{eq:incorrect-cocycle}
 \beta_p(n)=\beta_p(n+1)+\beta_p\bigl(n/(n+1)\bigr),
\end{equation}
where arguments are reduced modulo the prime $p$. The same statement
appears in the author-hosted preprint checked for this report
\cite{RZ}. For the exterior symbol defined in \eqref{eq:beta}, this
assertion is false. The issue is not a decimal discrepancy or an
unresolved branch convention: the obstruction is a nonzero rationalized
exterior symbol.

\begin{proposition}[Prime-seven counterexample]\label{prop:counterexample}
At $p=7$ and $n=2$, the defect in \eqref{eq:incorrect-cocycle} is
$3\beta_7(2)\ne0$.
\end{proposition}
\begin{proof}
Modulo $7$ one has $2/3=3$, and $3\equiv-2^{-1}$. The valid
symmetries \eqref{eq:symmetries} therefore give
$\beta_7(3)=-\beta_7(2)$. Equation
\eqref{eq:incorrect-cocycle} would say
$\beta_7(2)=-2\beta_7(2)$.

Here is an independent witness that $\beta_7(2)\ne0$, not relying on
the full kernel theorem. Write
\[
 a=\cl{1-\zeta_7},\quad b=\cl{1-\zeta_7^2},\quad
 c=\cl{1-\zeta_7^3}.
\]
Pairing $k$ and $-k$ in the definition gives
\begin{equation}\label{eq:beta7-wedge}
 \beta_7(2)=2(b\wedge a+c\wedge b+a\wedge c)
          =-2(b-a)\wedge(c-a).
\end{equation}
Let
\[
 A=\log(2\sin(\pi/7)),\quad B=\log(2\sin(2\pi/7)),\quad
 C=\log(2\sin(3\pi/7)).
\]
Since $0<\pi/7<2\pi/7<3\pi/7<\pi/2$, one has $A<B<C$.
Take logarithmic absolute values at the two embeddings indexed by
$1$ and $2$. The vectors corresponding to $b-a$ and $c-a$ are
respectively
\[
 (B-A,C-B),\qquad(C-A,A-B).
\]
Their determinant is
\begin{equation}\label{eq:det-witness}
 -(B-A)^2-(C-B)(C-A)<0.
\end{equation}
Thus their exterior product has a nonzero image, proving
\eqref{eq:beta7-wedge} nonzero. The defect is consequently
$3\beta_7(2)$, as claimed.
\end{proof}

\subsection{Precise correction and unaffected statements}

For the definition in \eqref{eq:beta}, the universal finite-field
three-term identity \eqref{eq:incorrect-cocycle} should be removed.
Any intended replacement involving an altered module, action or quotient
would need a new statement and proof. It is not justified by applying
rational Herglotz three-term relations directly to finite-field residues.

The sign and inverse symmetries remain true. The rational boundary
formula \eqref{eq:boundary-F} also remains true: it follows directly
from finite products, as proved in Proposition~\ref{prop:boundary}.
The current ProveIt manuscript already includes such a direct derivation
and explicitly separates it from a finite-field cocycle claim. The
new classification therefore does not inherit the erroneous premise.

This correction is deliberately version-specific. It identifies the
statement in the preprint pages actually inspected; it does not claim
that all editions of the journal article, all subsequent corrections,
or every cocycle construction in the paper have been checked. The
analytic Hecke relation and the finite dilogarithm evaluations used in
Section~\ref{sec:identities} do not use
\eqref{eq:incorrect-cocycle} in their displayed proofs.

\subsection{Editorial changes recommended for the manuscript}

The next integration can replace the open general-$J(p/q)$ formal
classification question with Theorem~\ref{thm:classification}, while
retaining the separate numerical period question. The phrase
``modulo $2e$'' should be preserved in every summary of the mixed-parity
criterion; the pair $(8,3)$ is a useful regression test.

The analytic identity \eqref{eq:neighbor} and the evaluation
\eqref{eq:J35} can be placed immediately after the existing $J(2/q)$
section. The finite-field caveat belongs beside the direct boundary
calculation. The theorem ledger should mark the existing kernel theorem,
the $K$-theory facts and the Hecke relation as prior inputs, not new
results. No earlier Gaussian or $S_p$ conjecture should have its status
changed as a side effect of this integration.

\section{Reproducibility, exact tests and numerical diagnostics}
\label{sec:verification}

The package separates three forms of evidence: the all-parameter ordinary
proofs in this article, finite exact regressions for the implementation,
and independent high-precision numerical comparisons. None is described
as proof-assistant verification. The numeric script is not an interval
arithmetic implementation and does not claim rigorous error enclosures.

\subsection{Exact finite checks}

The script \nolinkurl{code/verify_exact.py} uses integer dictionaries for
group-ring and exterior-algebra calculations, integer modular arithmetic
for classification, and SymPy only for exact polynomial and rational
simplifications. The frozen run uses a bound of $500$ and records:
\begin{center}
\begin{tabular}{p{0.71\linewidth}r}
\toprule
Check & Cases\\
\midrule
First-difference criterion over odd conductors & $50\,700$\\
Second-difference criterion over odd conductors & $50\,700$\\
New dyadic layer four-term group-ring witness & $50\,828$\\
Odd conductor doubling in a free exterior algebra & $2\,106$\\
Coprime unordered pairs checked against the recurrence generator & $76\,115$\\
Height-count regressions at every bound from $1$ to $500$ & $500$\\
Leading recurrence-polynomial coefficient checks & $62$\\
\bottomrule
\end{tabular}
\end{center}

The first two scans use odd conductors at most $500$. The new-layer
scan uses even $m\le500$ and top conductor $2m$. The free exterior
checks use odd conductors at most $101$, using the odd lift of a unit
residue in the conductor-doubling formula. They verify the expanded
identity before any further multiplicative unit relations are imposed.
The polynomial tests run through degree $32$ in each sign.

The brute-force arithmetic scan finds $781$ unordered logarithmic pairs
of height at most $500$. It finds exactly the three unordered
non-logarithmic rational-dilogarithmic pairs $(1,3)$, $(1,5)$ and
$(3,5)$. Independent recurrence generation produces exactly the same
logarithmic set. The script also checks the prime-seven defect and the
rational substitutions in both analytic five-term proofs. Frozen
regressions include the failure at $(8,3)$ and successful cases
$(12,5)$ and $(4,15)$.

These are meaningful exact cross-checks of implementations: the
recurrence generator does not call the modular classifier, and the
four-term group-ring calculation does not call a logarithmic numerical
rank estimator. Nevertheless their finite ranges cannot establish the
infinite classifications on their own.

\subsection{Independent defining-integral comparisons}

For reduced $p/q>0$, the numerical implementation evaluates the defining
integral after $t=u^q$:
\begin{equation}\label{eq:numeric-integral}
 J(p/q)=\int_0^1
       \frac{q u^{q-1}\log(1+u^p)}{1+u^q}\dd u.
\end{equation}
All exponents are nonnegative integers, making this a smooth integrand
at zero. The right sides are evaluated independently using log-sine
sums or classical dilogarithms, not by another invocation of the same
Herglotz identity.

The frozen script \nolinkurl{code/verify_numeric.py} performs $38$
comparisons at $90$ decimal digits and repeats them at $140$ digits.
There are $23$ neighboring cases ($1\le n\le20$, then $30$, $50$ and
$100$), the five values $J(1/3),J(1/5),J(3/5),J(4/3),J(4/5)$, the
reciprocal value $J(5/3)$, the exceptional cancellation identity, and
eight further reciprocal pairs. Every residual passes the diagnostic
threshold $10^{-d+15}$ at working precision $d$.

For orientation, the independently compared value is
\[
 J(3/5)=0.30458660277348562359158402248120510653743219820653\ldots.
\]
The output also records a negative numerical value for the determinant
in \eqref{eq:det-witness}. Its rigorous sign is established by the
inequality proof, not by this numerical check.

\subsection{Files, environment and replay}

The recorded environment is Python~3.13.5, SymPy~1.14.0 and
mpmath~1.3.0. The Python dependencies are pinned in
\nolinkurl{requirements.txt}. The principal commands, from the package
root, are
\begin{verbatim}
python -m pip install -r requirements.txt
python code/verify_exact.py --bound 500
python code/verify_numeric.py --dps 90 140
python code/classify.py 3 5
latexmk -pdf -interaction=nonstopmode -halt-on-error article.tex
\end{verbatim}
An existing TeX installation with the packages listed in
\nolinkurl{article.tex} is needed for the last command. The scripts do
not contact the network after their dependencies have been installed.

The JSON files \nolinkurl{data/exact_results.json} and
\nolinkurl{data/numeric_results.json} preserve the actual runs.
\nolinkurl{data/logarithmic_pairs.csv} records each pair with its
recurrence witnesses, while \nolinkurl{data/counts.csv} contains exact
larger-height counts. The theorem ledger distinguishes proved results,
known inputs and open questions. The integration directory contains an
additive manuscript fragment and a proposed bibliography entry; it does
not overwrite the existing manuscript or assert that a repository commit
has been made.

\section{Further research questions}
\label{sec:future}

The formal classification closes one stated problem, but it also makes
the next tasks more precise. The questions below are not used as
assumptions in any theorem of this article.

\subsection{Short complete certificates for every recurrence pair}

\begin{question}\label{question:certificates}
Given adjacent terms of an even-coefficient recurrence
\eqref{eq:minus-rec} or \eqref{eq:plus-rec}, construct the full
logarithmic correction for $J(a/b)$ with a replayable branch certificate.
Can the number of Hecke and five-term law instances be bounded by a
constant times $\log\max(a,b)$ after consecutive pairs are handled by
Theorem~\ref{thm:neighbor}?
\end{question}

The classification proves existence by rationalized Bloch descent. It
does not automatically make that descent an efficient symbolic
algorithm. A useful implementation should keep the finite log-sine sums
or cyclotomic products in a compressed representation, and state
separately the number of functional-equation steps, the expanded output
size and the bit complexity of algebraic-number operations. An
$O(\log B)$ descent path does not by itself imply an $O(\log B)$
expanded expression.

The neighboring theorem gives a complete base family. The next test
families are the coefficient-$2$ plus sequence
$1,2,5,12,29,\ldots$ and the coefficient-$4$ minus sequence
$1,4,15,56,\ldots$. Their exact classification means a failed numerical
search should now be treated as a defect of the proposed basis or the
certificate construction, not as evidence of formal non-reducibility.

\subsection{Other prime filters}

\begin{question}\label{question:prime-filter}
For a prime $r>2$, classify the formal boundary of
$F(rx)-2F(x)+F(x/r)$ at all positive rational $x$, including the cases
where the $r$-adic valuations of numerator and denominator change the
cyclotomic fields.
\end{question}

At $r=2$, the new layer has an order-two kernel and the four-term witness
in \eqref{eq:four-terms} is decisive. For odd $r$, the analogous new
character space is larger and its sparse group-algebra annihilator must
be analyzed afresh. Simply replacing $2$ by $r$ in
\eqref{eq:mixed-class} has not been justified. A first objective is an
exact conductor-$r$ distribution identity and a precise old-field
intersection theorem.

\subsection{Integral structure and torsion}

\begin{question}\label{question:integral}
What integral lattices underlie the symbol spaces and their dyadic
old/new decompositions? Which denominators and torsion classes are
necessary in explicit five-term certificates?
\end{question}

Rationalization is essential to the clean criterion
\eqref{eq:zero-detection}. It deliberately suppresses the torsion
structure and branch constants that an integral or flattened Bloch-group
implementation must retain. Smith normal forms in a finite
presentation could help locate certificate denominators, but a finite
matrix must not be mistaken for the entire integral pre-Bloch group.

\subsection{Sharper exact counting and uniform depth}

\begin{question}\label{question:counting}
Can the fixed-$R$ expansion \eqref{eq:count-expansion} be made effective
uniformly for $R$ growing with $B$, and can the remaining floor-function
oscillations be isolated by a short exact formula?
\end{question}

The proof gives each fixed fractional-power term and a deliberately
coarse tail. A refined treatment should separate the nearly degenerate
coefficient-$2$ minus family from the geometrically growing families,
and quantify the inversion of the recurrence polynomials uniformly in
their degree. This is a counting problem about the proved formal set,
not a numerical-independence problem.

\subsection{Period evaluation versus formal obstruction}

\begin{question}\label{question:periods}
For which explicitly defined subspaces generated by rational Herglotz
values is numerical period evaluation injective modulo products of
algebraic logarithms and $\pi^2$? Is there an arithmetic statement
strong enough to turn selected nonzero symbols into genuine numerical
non-reduction results?
\end{question}

This question is not answered by the present theorem. The six-extra
classification is a theorem about the rationalized five-term calculus.
Promoting it to an unconditional assertion about all real-number
identities would require an additional theorem of a different kind.
Any conjecture in that direction should first specify the coefficient
field, the permitted logarithmic products and the precise comparison
algebra.

\subsection{A corrected cocycle formulation and colored extensions}

\begin{question}\label{question:cocycle}
Which module and group action, if any, give a valid finite-field cocycle
related to the symbols $\beta_p(a)$ while accommodating the explicit
prime-seven defect? How does that formulation interact with the
Herglotz--Zagier--Novikov parameters of \cite{CK}?
\end{question}

A proposed identity must preserve the nonzero wedge detected in
\eqref{eq:det-witness}, or explicitly state a quotient that kills it.
The rational boundary formula is available independently, so such a
reformulation need not be forced into the proof of the classification.
Colored parameters may require replacing the real positive Gram matrix
by a character-sensitive nondegeneracy argument. This is a distinct
extension, not a consequence of the dyadic proof by notation change.

\subsection{Formal verification targets}

A realistic first formalization would cover the finite group algebra,
normalized norm identities, parity-sensitive boundary expansion and
Vieta descent. The analytic component could then formalize the Rogers
identity on $(0,1)$ and replay the two substitutions in
Section~\ref{sec:identities}. The classical $K$-theory descent and
Dirichlet nonvanishing inputs should remain explicit interfaces until
they are themselves available in the chosen proof assistant.

\section*{Conclusion}
\addcontentsline{toc}{section}{Conclusion}

The all-rational formal $J$ problem is governed by a dyadic old/new
separation, not by a search over increasingly large baskets of constants.
That separation yields a complete modular test, a complete quadratic
recurrence parametrization and a hierarchy of exact height asymptotics.
The analytic Hecke relation then produces a full infinite identity and
a compact exceptional dilogarithm evaluation. The finite-field audit
illustrates why direct symbol calculations, branch control and the
distinction between formal and numerical statements must remain visible
throughout the continuation.

\clearpage
\phantomsection
\addcontentsline{toc}{section}{References}
\begin{thebibliography}{9}

\bibitem{Repo}
ProveIt contributors,
\emph{Polylogarithms and their Arithmetic Bridges},
consolidated manuscript, source snapshot
\texttt{cc34f73596336f2466d9754cb0f3635bd2bedade}.
See the chapter files \nolinkurl{09-herglotz.tex} and
\nolinkurl{10-discovery.tex}, and the report
\nolinkurl{herglotz-cyclotomic-obstructions}.
\url{https://github.com/VladimirReshetnikov/ProveIt/tree/cc34f73596336f2466d9754cb0f3635bd2bedade/Analysis/Polylogarithms/docs/manuscript}.

\bibitem{RZ}
D. Radchenko and D. Zagier,
\emph{Arithmetic properties of the Herglotz function},
J. reine angew. Math. \textbf{797} (2023), 229--253.
The formula numbering and the correction in this article refer to the
reviewed preprint \texttt{arXiv:2012.15805v1} (2020), especially
(15), (22a), (23), (24), Theorem 3 and Section 7.2.
\url{https://arxiv.org/abs/2012.15805v1}.
Author-hosted preprint:
\url{https://people.mpim-bonn.mpg.de/zagier/files/preprints/HerglotzFunction.pdf}.

\bibitem{CK}
Y. Choie and R. Kumar,
\emph{Arithmetic properties of the Herglotz--Zagier--Novikov function},
\texttt{arXiv:2309.10634v1} (2023).
\url{https://arxiv.org/abs/2309.10634v1}.

\bibitem{ZagierDilog}
D. Zagier,
\emph{The dilogarithm function}, in
P. Cartier et al. (eds.), \emph{Frontiers in Number Theory, Physics, and
Geometry II}, Springer, 2007, pp. 3--65.
\url{https://doi.org/10.1007/978-3-540-30308-4_1}.
Author's full text:
\url{https://people.mpim-bonn.mpg.de/zagier/files/doi/10.1007/978-3-540-30308-4_1/fulltext.pdf}.

\bibitem{CGZ}
F. Calegari, S. Garoufalidis and D. Zagier,
\emph{Bloch groups, algebraic K-theory, units, and Nahm's conjecture},
Ann. Sci. \'Ec. Norm. Sup\'er. (4) \textbf{56} (2023), 383--426.
\url{https://doi.org/10.24033/asens.2537}.
Author's full text: \url{https://math.uchicago.edu/~fcale/papers/CGZ.pdf}.

\bibitem{DLMF}
NIST Digital Library of Mathematical Functions,
\emph{Section 25.15: Dirichlet $L$-functions},
especially Equation 25.15.9, $L(1,\chi)\ne0$ for nonprincipal $\chi$.
\url{https://dlmf.nist.gov/25.15}.

\end{thebibliography}

\end{document}
